\chapter{Referencial Teórico}
\label{cap:referencial}

Este capítulo apresenta os fundamentos que sustentam o trabalho: a Mineração de Dados Educacionais e o problema da opacidade algorítmica, o algoritmo XGBoost, as definições formais de Justiça Algorítmica, a Inteligência Artificial Explicável baseada em valores de Shapley e a metodologia de \textit{Design Science Research}.

% ---------------------------------------------------------------------------
\section{Mineração de Dados Educacionais e Predição de Desempenho}
\label{sec:rt_mde}

A Mineração de Dados Educacionais (MDE) aplica técnicas de estatística e aprendizado de máquina a dados produzidos em contextos de ensino, com o objetivo de compreender processos de aprendizagem e apoiar decisões pedagógicas e administrativas \cite{romero2010educational}. Entre suas aplicações mais difundidas estão a predição de desempenho acadêmico, a identificação precoce de risco de evasão e o apoio a processos seletivos.

No Brasil, os microdados do ENEM constituem uma das bases mais utilizadas para esse fim, por combinarem as notas de milhões de participantes com um questionário socioeconômico detalhado \cite{inep2022microdados}. Essa riqueza, contudo, traz um risco: as variáveis do questionário (renda familiar, escolaridade dos pais, raça/cor, acesso a tecnologia) registram desigualdades estruturais da sociedade brasileira. Modelos treinados sobre esses dados tendem a incorporá-las como regularidades estatísticas. \citeonline{baker2022algorithmic} mostram que o viés algorítmico em sistemas educacionais é documentado para diversos grupos demográficos e pode surgir em qualquer etapa do ciclo de vida do modelo, da coleta dos dados à interpretação das predições.

% ---------------------------------------------------------------------------
\section{Opacidade Algorítmica e o Problema da Caixa-Preta}
\label{sec:rt_caixa_preta}

Modelos de alto desempenho, como \textit{ensembles} de árvores e redes neurais, combinam centenas ou milhares de componentes cujas interações não são diretamente inteligíveis para seres humanos. Essa característica é conhecida como problema da \textit{caixa-preta} \cite{adadi2018peeking}. A opacidade dificulta a responsabilização, impede a contestação de decisões e oculta padrões discriminatórios. \citeonline{rudin2019stop} argumenta que, em decisões de alto impacto, modelos intrinsecamente interpretáveis deveriam ser preferidos. Quando modelos complexos são empregados, técnicas de explicação \textit{post-hoc} tornam-se condição mínima de transparência \cite{arrieta2020explainable}.

A questão ganha contornos jurídicos no Brasil. O Art.~20 da Lei Geral de Proteção de Dados assegura ao titular o direito de solicitar a revisão de decisões tomadas unicamente com base em tratamento automatizado e de obter informações sobre os critérios utilizados \cite{brasil2018lgpd}. O Projeto de Lei nº 2.338/2023 tende a classificar sistemas de IA aplicados a processos seletivos educacionais como de alto risco \cite{brasil2023pl2338}.

% ---------------------------------------------------------------------------
\section{\textit{Gradient Boosting} e o Algoritmo XGBoost}
\label{sec:rt_xgboost}

O \textit{gradient boosting} constrói um modelo preditivo como soma de modelos fracos, tipicamente árvores de decisão rasas, ajustados sequencialmente para corrigir os erros dos anteriores \cite{friedman2001greedy, hastie2009elements}. Para uma instância $x_i$, a predição após $K$ árvores é:
\begin{equation}
\hat{y}_i = \sum_{k=1}^{K} f_k(x_i), \qquad f_k \in \mathcal{F},
\label{eq:rt_boosting}
\end{equation}
em que $\mathcal{F}$ é o espaço das árvores de regressão. Em classificação binária, $\hat{y}_i$ é um escore na escala log-odds, convertido em probabilidade pela função logística $\sigma(z) = 1/(1+e^{-z})$.

O XGBoost \cite{chen2016xgboost} formaliza o treinamento como a minimização de um objetivo regularizado:
\begin{equation}
\mathcal{L} = \sum_{i=1}^{n} \ell(y_i, \hat{y}_i) + \sum_{k=1}^{K} \Omega(f_k), \qquad \Omega(f) = \gamma T + \frac{1}{2}\lambda \lVert w \rVert^2,
\label{eq:rt_xgb_obj}
\end{equation}
em que $\ell$ é a função de perda (log-loss, no caso binário), $T$ é o número de folhas da árvore, $w$ é o vetor de pesos das folhas e $\gamma$ e $\lambda$ são hiperparâmetros de regularização. A cada iteração, o algoritmo aproxima a perda por uma expansão de Taylor de segunda ordem, o que permite calcular analiticamente o peso ótimo de cada folha e o ganho de cada divisão. Somados a otimizações de sistema (paralelização e tratamento de esparsidade), esses elementos explicam o desempenho do XGBoost em dados tabulares, seu domínio característico. A estrutura em árvores também viabiliza o cálculo exato e eficiente de explicações (Seção~\ref{sec:rt_xai}).

% ---------------------------------------------------------------------------
\section{Justiça Algorítmica}
\label{sec:rt_fairness}

\subsection{Origens do Viés}

A Justiça Algorítmica (\textit{Algorithmic Fairness}) estuda como sistemas automatizados podem produzir resultados sistematicamente desfavoráveis a grupos definidos por atributos sensíveis, como raça, gênero ou renda, e como medir e mitigar esse efeito \cite{barocas2023fairness, mehrabi2021survey}. Entre as fontes de viés destacam-se:
\begin{itemize}
    \item o \textbf{viés histórico}, quando os dados refletem desigualdades preexistentes na sociedade, mesmo que corretamente medidos;
    \item o \textbf{viés de representação}, quando grupos minoritários são sub-representados, o que torna as estimativas sobre eles menos precisas;
    \item o \textbf{viés de rótulo}, quando a variável-alvo é, ela própria, produto de processos discriminatórios.
\end{itemize}

Uma resposta ingênua seria simplesmente remover o atributo sensível do modelo (\textit{fairness through unawareness}). \citeonline{dwork2012fairness} demonstram que essa estratégia é insuficiente, pois outras variáveis correlacionadas ao atributo (\textit{proxies}) permitem reconstruí-lo implicitamente. Por isso, a auditoria precisa medir os \textbf{resultados} do modelo por grupo, e não apenas inspecionar suas entradas.

\subsection{Critérios de Justiça de Grupo}

Sejam $\hat{Y} \in \{0,1\}$ a predição, $Y \in \{0,1\}$ o rótulo real e $G$ o atributo sensível. A literatura consolidou três famílias principais de critérios de justiça de grupo \cite{barocas2023fairness}.

A \textbf{Paridade Demográfica} (independência) exige que a probabilidade de predição positiva não dependa do grupo:
\begin{equation}
P(\hat{Y} = 1 \mid G = a) = P(\hat{Y} = 1 \mid G = b), \quad \forall\, a, b.
\label{eq:rt_dp}
\end{equation}
Sua forma de razão é o \textbf{Impacto Disparate} (DI), que compara a taxa de seleção de um grupo $a$ à de um grupo de referência $b$:
\begin{equation}
\mathrm{DI} = \frac{P(\hat{Y} = 1 \mid G = a)}{P(\hat{Y} = 1 \mid G = b)}.
\label{eq:rt_di}
\end{equation}
A métrica tem origem no direito antidiscriminatório norte-americano. As diretrizes da \textit{Equal Employment Opportunity Commission} estabelecem a \textbf{regra dos quatro quintos}: uma taxa de seleção inferior a 80\% da taxa do grupo mais favorecido constitui evidência de impacto adverso \cite{eeoc1978uniform}. \citeonline{feldman2015certifying} formalizaram essa regra no contexto do aprendizado de máquina, definindo como discriminatório um classificador com $\mathrm{DI} < \tau$, com $\tau = 0{,}8$.

A Paridade Demográfica ignora os rótulos reais. Se as taxas-base diferem entre os grupos, mesmo um classificador perfeito a violaria. Para considerar a correção das predições, \citeonline{hardt2016equality} propuseram a \textbf{Igualdade de Oportunidades} (\textit{Equal Opportunity}), que exige taxas de verdadeiros positivos iguais:
\begin{equation}
P(\hat{Y} = 1 \mid Y = 1, G = a) = P(\hat{Y} = 1 \mid Y = 1, G = b),
\label{eq:rt_eo}
\end{equation}
e sua versão mais forte, as \textbf{Chances Equalizadas} (\textit{Equalized Odds}), que exige também a igualdade das taxas de falsos positivos:
\begin{equation}
P(\hat{Y} = 1 \mid Y = y, G = a) = P(\hat{Y} = 1 \mid Y = y, G = b), \quad y \in \{0, 1\}.
\label{eq:rt_eodds}
\end{equation}
No contexto educacional, a Igualdade de Oportunidades tem interpretação direta: candidatos com desempenho efetivamente alto devem ter a mesma chance de serem reconhecidos pelo modelo, qualquer que seja o seu grupo.

\subsection{Incompatibilidade entre Critérios}

Os critérios acima não podem, em geral, ser satisfeitos simultaneamente. \citeonline{kleinberg2017inherent} e \citeonline{chouldechova2017fair} demonstraram que, quando as taxas-base diferem entre os grupos, é matematicamente impossível equalizar ao mesmo tempo a calibração e as taxas de erro, exceto em casos degenerados. Consequentemente, uma auditoria não deve se apoiar em um único indicador: a escolha das métricas é uma decisão normativa, e reportar critérios complementares, como o DI e a TPR, oferece um diagnóstico mais completo. Bibliotecas como Fairlearn \cite{bird2020fairlearn} e AI Fairness 360 \cite{bellamy2018aif360} implementam esses critérios, mas exigem integração programática especializada.

% ---------------------------------------------------------------------------
\section{Inteligência Artificial Explicável}
\label{sec:rt_xai}

\subsection{Taxonomia}

A Inteligência Artificial Explicável (XAI) reúne métodos que tornam compreensíveis as decisões de modelos de aprendizado de máquina \cite{arrieta2020explainable, adadi2018peeking}. Esses métodos podem ser classificados em três eixos:
\begin{itemize}
    \item \textbf{intrínsecos} (o modelo é interpretável por construção) ou \textbf{\textit{post-hoc}} (aplicados a um modelo já treinado);
    \item \textbf{globais} (explicam o comportamento geral do modelo) ou \textbf{locais} (explicam uma predição individual);
    \item \textbf{agnósticos ao modelo} ou \textbf{específicos} de uma família de modelos.
\end{itemize}
O método SHAP, adotado neste trabalho, é \textit{post-hoc}, produz explicações locais agregáveis em globais e possui uma versão específica e exata para modelos baseados em árvores.

\subsection{Valores de Shapley}

Os valores de Shapley provêm da teoria dos jogos cooperativos e resolvem o problema de repartir o ganho de uma coalizão entre seus jogadores \cite{shapley1953value}. Considere um conjunto $F$ de $M$ jogadores e uma função característica $v: 2^F \rightarrow \mathbb{R}$ que atribui um valor a cada coalizão $S \subseteq F$. A contribuição do jogador $i$ é a média ponderada de suas contribuições marginais sobre todas as coalizões possíveis:
\begin{equation}
\phi_i(v) = \sum_{S \subseteq F \setminus \{i\}} \frac{|S|!\,(M - |S| - 1)!}{M!}\, \bigl[ v(S \cup \{i\}) - v(S) \bigr].
\label{eq:rt_shapley}
\end{equation}

O valor de Shapley é a única regra de repartição que satisfaz simultaneamente quatro axiomas:
\begin{enumerate}
    \item \textbf{eficiência}: as contribuições somam o valor total, $\sum_i \phi_i = v(F) - v(\emptyset)$;
    \item \textbf{simetria}: jogadores com contribuições marginais idênticas recebem valores iguais;
    \item \textbf{jogador nulo}: um jogador que não altera nenhuma coalizão recebe valor zero;
    \item \textbf{aditividade}: o valor de um jogo composto é a soma dos valores dos jogos componentes.
\end{enumerate}

\subsection{SHAP e TreeSHAP}

\citeonline{lundberg2017unified} transpuseram esse arcabouço para a explicação de modelos. Na formulação SHAP (\textit{SHapley Additive exPlanations}), os atributos são os jogadores e a predição do modelo é o ganho. A explicação assume a forma de um modelo aditivo de atribuição:
\begin{equation}
g(z') = \phi_0 + \sum_{i=1}^{M} \phi_i z'_i, \qquad z' \in \{0,1\}^M,
\label{eq:rt_shap}
\end{equation}
em que $z'_i$ indica a presença do atributo $i$ e $\phi_0$ é a predição média do modelo. Os autores demonstram que os valores de Shapley constituem a única solução dessa classe que satisfaz as propriedades de precisão local, ausência (\textit{missingness}) e consistência. A precisão local garante que a explicação reconstrói exatamente a predição: $f(x) = \phi_0 + \sum_i \phi_i$.

O cálculo direto da Equação~\eqref{eq:rt_shapley} tem custo exponencial em $M$. Para modelos baseados em árvores, o algoritmo TreeSHAP calcula os valores exatos em tempo polinomial, $O(TLD^2)$, sendo $T$ o número de árvores, $L$ o número máximo de folhas e $D$ a profundidade máxima \cite{lundberg2020local}. Isso torna viável explicar todas as instâncias de um conjunto de teste. Agregados pela média dos valores absolutos, os valores locais fornecem uma medida de importância global consistente.

% ---------------------------------------------------------------------------
\section{\textit{Design Science Research}}
\label{sec:rt_dsr}

A \textit{Design Science Research} (DSR) é um paradigma de pesquisa orientado à criação e à avaliação de artefatos (construtos, modelos, métodos ou instâncias) que resolvem problemas relevantes \cite{hevner2004design}. Em contraste com a pesquisa comportamental, que busca explicar fenômenos existentes, a DSR produz conhecimento ao projetar algo novo e demonstrar rigorosamente sua utilidade. \citeonline{hevner2004design} enfatizam que a avaliação do artefato, por métodos analíticos, experimentais ou observacionais, é parte constitutiva da contribuição científica.

\citeonline{peffers2007design} sistematizaram o paradigma na \textit{Design Science Research Methodology}, composta por seis atividades:
\begin{enumerate}
    \item identificação do problema e motivação;
    \item definição dos objetivos da solução;
    \item projeto e desenvolvimento;
    \item demonstração;
    \item avaliação;
    \item comunicação.
\end{enumerate}
O processo é iterativo, e os resultados da avaliação podem realimentar o projeto. A aplicação dessa metodologia ao presente trabalho é detalhada no Capítulo~\ref{cap:metodologia}.

% ---------------------------------------------------------------------------
\section{Considerações do Capítulo}

Os fundamentos apresentados delineiam a lacuna que o trabalho aborda. Modelos como o XGBoost oferecem alto desempenho em dados tabulares, mas são opacos. Os critérios de justiça de grupo permitem quantificar disparidades, mas nenhum deles é suficiente isoladamente. O SHAP permite explicar as decisões com garantias teóricas, e a DSR fornece o arcabouço para construir e avaliar rigorosamente um artefato que integre esses elementos.
